% Worked GAIA example: planner JSON -> procedure DAG -> C(Q)
\subsection{Worked Example: GAIA Query to Procedure DAG and Complexity Vector}
\label{app:gaia-worked-example}

We walk through a real held-out GAIA transfer query (eval sample
\texttt{f46b4380-207e-4434-820b-f32ce04ae2a4}, Level~2) from decomposition through graph repair to
the five-dimensional complexity vector $\mathbf{z}(q)$ consumed by the router.

\paragraph{Query.}
\begin{quote}\small
It is 1999. Before you party like it is 1999, please assist me in settling a bet.
Fiona Apple and Paula Cole released albums prior to 1999. Of these albums, which didn't receive
a letter grade from Robert Christgau? Provide your answer as a comma delimited list of album titles,
sorted alphabetically.
\end{quote}

\paragraph{Planner output (JSON).}
The procedure planner returns \texttt{final\_constraint=phrase} and four subtasks with
multiple \texttt{depends\_on} edges (no graph repairs on this instance):

\begin{table}[h]
  \centering
  \footnotesize
  \caption{Planner subtasks for the GAIA album-grade query (abridged).}
  \label{tab:gaia-subtasks}
  \begin{tabular}{@{}lllll@{}}
    \toprule
    id & step\_type & tool & depends\_on & goal (abridged) \\
    \midrule
    t1 & retrieve & web\_search & [] & Fiona Apple albums before 1999 \\
    t2 & retrieve & web\_search & [] & Paula Cole albums before 1999 \\
    t3 & retrieve & web\_search & [t1, t2] & Christgau letter grades for albums \\
    t4 & compute & none & [t1, t2, t3] & Albums without letter grade (sink) \\
    \bottomrule
  \end{tabular}
\end{table}

\paragraph{Graph statistics and $\mathbf{z}(q)$.}
QCE computes topology and plan-semantics statistics on $G_q$ (e.g., \texttt{max\_depth}$=3$,
\texttt{n\_nodes}$=4$, multiple merge edges, \texttt{n\_merge\_nodes}${=}2$),
normalizes each statistic with train-fitted caps (\texttt{models/qce\_graph/train\_norm.json}),
and aggregates normalized terms into five groups
(Table~\ref{tab:complexity-dims}; implementation in \texttt{qce/graph.py}).

\begin{table}[h]
  \centering
  \footnotesize
  \caption{Complexity vector for the GAIA album-grade example (train norm applied).}
  \label{tab:gaia-cq}
  \begin{tabular}{@{}lrl@{}}
    \toprule
    Dimension & Value & Main signals (this query) \\
    \midrule
    $z_1$ structural & 0.883 & 4-node DAG, parallel roots, multi-parent sink \\
    $z_2$ reasoning & 0.508 & compute sink after merged retrieval \\
    $z_3$ evidence & 0.719 & three web retrieval hops \\
    $z_4$ tool & 0.486 & repeated \texttt{web\_search}, no attachment \\
    $z_5$ coordination & 0.143 & clean plan; non-trivial merge structure \\
    \bottomrule
  \end{tabular}
\end{table}

\paragraph{Procedure DAG.}
QCE links virtual \texttt{\_\_start\_\_} to both root retrieves, connects each
\texttt{depends\_on} prerequisite, and attaches \texttt{\_\_end\_\_} after the sink
(Figure~\ref{fig:gaia-dag-example}).
The graph is not a simple chain: \texttt{t3} merges \texttt{t1} and \texttt{t2}, and
\texttt{t4} depends on all three predecessors.

% GAIA worked example — top-to-bottom procedure DAG (eval f46b4380)
\begin{figure*}[t]
  \centering
  \resizebox{0.98\textwidth}{!}{%
  \begin{tikzpicture}[
    font=\small\sffamily,
    >=Stealth,
    panel/.style={
      draw=gray!45, rounded corners=4pt, fill=gray!4,
      inner sep=6pt, align=left, font=\scriptsize\sffamily,
    },
    st/.style={
      draw=#1!65!black, fill=#1!12, rounded corners=6pt,
      minimum width=2.0cm, minimum height=1.15cm, align=center,
      line width=0.95pt, inner sep=5pt,
    },
    virtnode/.style={
      draw=gray!50, fill=white, line width=0.75pt,
      rounded corners=4pt, font=\scriptsize\sffamily\ttfamily,
      inner xsep=6pt, inner ysep=3pt,
    },
    dep/.style={->, line width=1.1pt, draw=blue!65!black},
    depside/.style={->, line width=0.85pt, draw=blue!35},
    virtedge/.style={->, line width=0.8pt, draw=gray!55, dashed},
  ]
    % ===== Top metadata =====
    \node[panel, text width=3.85cm] (ctx) at (-4.85, 9.2) {%
      \textbf{Context}\\[2pt]
      1999 bet: Fiona Apple \& Paula Cole albums---which lacked a Christgau \emph{letter} grade? Comma-separated titles (A--Z).
    };
    \node[panel, text width=2.4cm, align=center] (qinfo) at (0, 9.2) {%
      \textbf{Query} $\mathbf{Q}$\\
      GAIA L2 \texttt{f46b4380\ldots}
    };
    \node[panel, text width=3.9cm] (edgebox) at (4.85, 9.2) {%
      \textbf{\texttt{depends\_on}}\\
      \texttt{t1[]}, \texttt{t2[]}; \texttt{t3[t1,t2]}; \texttt{t4[t1,t2,t3]}
    };

    % ===== DAG background (white, behind nodes) =====
    \fill[white, rounded corners=10pt] (-4.05, 0.65) rectangle (4.05, 8.35);
    \draw[gray!45, rounded corners=10pt, line width=0.75pt]
      (-4.05, 0.65) rectangle (4.05, 8.35);

    % ===== DAG nodes: start (top) -> end (bottom) =====
    \node[virtnode] (start) at (0, 8.15) {start};

    \node[st=blue] (t1) at (-2.75, 6.65) {%
      \textbf{t1}\\[2pt]
      {\scriptsize retrieve}\\
      {\scriptsize\texttt{web\_search}}};
    \node[st=blue] (t2) at (2.75, 6.65) {%
      \textbf{t2}\\[2pt]
      {\scriptsize retrieve}\\
      {\scriptsize\texttt{web\_search}}};

    \node[st=blue] (t3) at (0, 4.75) {%
      \textbf{t3}\\[2pt]
      {\scriptsize retrieve}\\
      {\scriptsize\texttt{web\_search}}};

    \node[st=violet] (t4) at (0, 2.65) {%
      \textbf{t4}\\[2pt]
      {\scriptsize compute}\\
      {\scriptsize sink}};

    \node[virtnode] (end) at (0, 1.15) {end};

    % ===== DAG edges (downward flow) =====
    \draw[virtedge] (start.south west) -- (t1.north);
    \draw[virtedge] (start.south east) -- (t2.north);
    \draw[dep] (t1.south) -- (t3.north west);
    \draw[dep] (t2.south) -- (t3.north east);
    \draw[dep] (t3.south) -- (t4.north);
    \draw[virtedge] (t4.south) -- (end.north);
    \draw[depside] (t1.south east) to[out=-58, in=168, looseness=0.85] (t4.north west);
    \draw[depside] (t2.south west) to[out=-122, in=12, looseness=0.85] (t4.north east);

    % ===== Bottom metadata =====
    \node[panel, text width=3.5cm] (leg) at (-3.15, -0.35) {%
      \textbf{Legend}\\
      \tikz{\draw[dep] (0,0) -- (0.55,0);} spine \quad
      \tikz{\draw[depside] (0,0) -- (0.55,0);} direct\\
      \tikz{\draw[virtedge] (0,0) -- (0.55,0);} virtual
    };
    \node[panel, text width=5.2cm, align=left] (zq) at (2.75, -0.35) {%
      \textbf{$\mathbf{z}(q)$} (train norm)\\[3pt]
      \begin{tabular}{@{}l@{\hspace{0.3em}}c@{\hspace{0.2em}}r@{}}
        Structural    & \textcolor{blue!55!black}{\rule{1.5cm}{0.12cm}} & 0.883 \\
        Reasoning     & \textcolor{blue!55!black}{\rule{0.92cm}{0.12cm}} & 0.508 \\
        Evidence      & \textcolor{blue!55!black}{\rule{1.22cm}{0.12cm}} & 0.719 \\
        Tool          & \textcolor{blue!55!black}{\rule{0.88cm}{0.12cm}} & 0.486 \\
        Coordination  & \textcolor{blue!55!black}{\rule{0.20cm}{0.12cm}} & 0.143 \\
      \end{tabular}
    };

  \end{tikzpicture}%
  }

  \caption{Worked GAIA example (eval \texttt{f46b4380-207e-4434-820b-f32ce04ae2a4}): procedure
    DAG $G_q$ (top \texttt{start} $\rightarrow$ bottom \texttt{end}) with parallel retrieves
    \texttt{t1}, \texttt{t2}, merge at \texttt{t3}, and fan-in sink \texttt{t4}
    (Table~\ref{tab:gaia-subtasks}). Curved edges are direct \texttt{depends\_on} links to the sink.
    Normalized $\mathbf{z}(q)$ in Table~\ref{tab:gaia-cq}.}
  \label{fig:gaia-dag-example}
\end{figure*}


The router then concatenates $\mathbf{z}(q)$ with $\psi(q)$ (PCA-16 embedding of the query
text) to form $\phi(q)$ for cost-soft HGBM routing.
This example illustrates how multi-hop GAIA queries with parallel evidence gathering and
merge nodes yield higher structural and evidence complexity than single-file linear plans.
